\begin{filecontents*}{underfish_report_refs.bib}
@standard{gost7322017,
  title        = {ГОСТ 7.32--2017. Система стандартов по информации, библиотечному и издательскому делу. Отчет о научно-исследовательской работе. Структура и правила оформления},
  organization = {Межгосударственный совет по стандартизации, метрологии и сертификации},
  year         = {2017},
  location     = {Москва}
}
@online{springboot,
  title   = {Spring Boot Reference Documentation},
  author  = {{VMware Tanzu}},
  year    = {2026},
  url     = {https://docs.spring.io/spring-boot/},
  urldate = {2026-06-02}
}
@online{springsecurity,
  title   = {Spring Security Reference},
  author  = {{VMware Tanzu}},
  year    = {2026},
  url     = {https://docs.spring.io/spring-security/reference/},
  urldate = {2026-06-02}
}
@online{jwt,
  title   = {JSON Web Token (JWT)},
  author  = {{IETF}},
  year    = {2015},
  url     = {https://www.rfc-editor.org/rfc/rfc7519},
  urldate = {2026-06-02}
}
@online{oauth2,
  title   = {The OAuth 2.0 Authorization Framework},
  author  = {{IETF}},
  year    = {2012},
  url     = {https://www.rfc-editor.org/rfc/rfc6749},
  urldate = {2026-06-02}
}
@online{docker,
  title   = {Docker Documentation},
  author  = {{Docker Inc.}},
  year    = {2026},
  url     = {https://docs.docker.com/},
  urldate = {2026-06-02}
}
@online{kubernetes,
  title   = {Kubernetes Documentation},
  author  = {{The Kubernetes Authors}},
  year    = {2026},
  url     = {https://kubernetes.io/docs/},
  urldate = {2026-06-02}
}
@online{postgresql,
  title   = {PostgreSQL Documentation},
  author  = {{The PostgreSQL Global Development Group}},
  year    = {2026},
  url     = {https://www.postgresql.org/docs/},
  urldate = {2026-06-02}
}
@online{redis,
  title   = {Redis Documentation},
  author  = {{Redis Ltd.}},
  year    = {2026},
  url     = {https://redis.io/docs/latest/},
  urldate = {2026-06-02}
}
@online{kafka,
  title   = {Apache Kafka Documentation},
  author  = {{The Apache Software Foundation}},
  year    = {2026},
  url     = {https://kafka.apache.org/documentation/},
  urldate = {2026-06-02}
}
@online{minio,
  title   = {MinIO Object Storage Documentation},
  author  = {{MinIO, Inc.}},
  year    = {2026},
  url     = {https://min.io/docs/},
  urldate = {2026-06-02}
}
@online{prometheus,
  title   = {Prometheus Documentation},
  author  = {{Prometheus Authors}},
  year    = {2026},
  url     = {https://prometheus.io/docs/},
  urldate = {2026-06-02}
}
@online{loki,
  title   = {Grafana Loki Documentation},
  author  = {{Grafana Labs}},
  year    = {2026},
  url     = {https://grafana.com/docs/loki/latest/},
  urldate = {2026-06-02}
}
@online{timepad,
  title   = {Timepad: сервис для организаторов событий},
  author  = {{Timepad}},
  year    = {2026},
  url     = {https://timepad.ru/},
  urldate = {2026-06-02}
}
@online{yandexafisha,
  title   = {Яндекс Афиша},
  author  = {{Яндекс}},
  year    = {2026},
  url     = {https://afisha.yandex.ru/},
  urldate = {2026-06-02}
}
\end{filecontents*}

\documentclass[14pt,a4paper]{extarticle}

% Базовые пакеты, требуемые в задании.
\usepackage[T2A]{fontenc}
\usepackage[utf8]{inputenc}
\usepackage[russian,english]{babel}
\usepackage{geometry}
\usepackage{graphics}
\usepackage{graphicx}
\usepackage[backend=biber,style=gost-numeric,sorting=none]{biblatex}

% Дополнительные пакеты оформления. Они не заменяют обязательные пакеты выше.
\usepackage{cmap}
\usepackage{csquotes}
\usepackage{array}
\usepackage{longtable}
\usepackage{enumitem}
\usepackage{titlesec}
\usepackage{tocloft}
\usepackage{caption}
\usepackage{float}
\usepackage{hyperref}

\addbibresource{underfish_report_refs.bib}
\graphicspath{{docs/}{docs/diagrams/}{diagrams/}{}}

\geometry{left=30mm,right=15mm,top=20mm,bottom=20mm}
\linespread{1.3}
\setlength{\parindent}{1.25cm}
\setlength{\parskip}{0pt}
\setlist{nosep,leftmargin=1.25cm}
\captionsetup{labelsep=endash,justification=centering}
\hypersetup{hidelinks}

\titleformat{\section}{\normalfont\bfseries\centering\MakeUppercase}{\thesection}{1em}{}
\titleformat{\subsection}{\normalfont\bfseries}{\thesubsection}{1em}{}
\titleformat{\subsubsection}{\normalfont\bfseries}{\thesubsubsection}{1em}{}

\renewcommand{\contentsname}{Содержание}
\renewcommand{\refname}{Список использованных источников}
\renewcommand{\figurename}{Рисунок}
\renewcommand{\tablename}{Таблица}

% Рисунки оформляются по ГОСТ: изображение по центру, подпись под рисунком.
% Файл изображения хранится в docs/diagrams или в каталоге проекта LaTeX.
% Пример: \reportfigure[0.9\textwidth]{uml_component_architecture.png}{UML-диаграмма компонентов}
\newcommand{\reportfigure}[3][0.95\textwidth]{%
\begin{figure}[H]
\centering
\includegraphics[width=#1,keepaspectratio]{#2}
\caption{#3}
\end{figure}
}

\begin{document}

\begin{titlepage}
\begin{center}
Министерство науки и высшего образования Российской Федерации\\
\vspace{0.5cm}
\textbf{Наименование образовательной организации}\\
\vspace{1.5cm}
\textbf{ОТЧЕТ}\\
по проектной/научно-исследовательской работе\\
\vspace{0.5cm}
\textbf{Платформа для анонсирования и поиска мероприятий Underfish}\\
\vspace{0.5cm}
\textbf{Backend, DevOps и тестирование}\\
\end{center}

\vfill

\begin{flushright}
Исполнитель: \underline{\hspace{6cm}}\\
Группа: \underline{\hspace{6.8cm}}\\
Руководитель: \underline{\hspace{5.7cm}}\\
\end{flushright}

\vfill

\begin{center}
Город -- 2026
\end{center}
\end{titlepage}

\tableofcontents
\newpage

\section*{Введение}
\addcontentsline{toc}{section}{Введение}

Современные сервисы поиска мероприятий решают задачу агрегации афиш, однако в большинстве случаев ориентированы на крупные коммерческие события. Для малых локальных сообществ, независимых организаторов, культурных инициатив, небольших мастерских и неформальных площадок такие решения оказываются недостаточно удобными: продвижение требует дополнительных затрат, публикация события конкурирует с большим количеством коммерческого контента, а пользователь не всегда может найти мероприятие рядом с домом или в конкретном районе.

Проект Underfish представляет собой гиперлокальную платформу для анонсирования и поиска мероприятий. Система предназначена для публикации концертов, встреч, мастер-классов, выставок, лекций и других событий, а также для поиска мероприятий с использованием карты, тегов, фильтров по датам, подписок на сообщества и персональных отметок участия.

В рамках данной работы рассматривается backend-часть платформы, инфраструктура разработки и развертывания, а также подход к тестированию. Frontend-часть и макеты Figma используются как источник пользовательских сценариев: по ним определялись требования к REST API, структуре профилей, карточкам мероприятий, поиску на карте и операциям организатора.

Разработка проекта выполнялась в два этапа. На первом этапе был спроектирован монолитный вариант системы, который позволил быстро проверить предметную область, основные сущности и сценарии работы пользователей. На втором этапе архитектура была декомпозирована на микросервисы, что позволило разделить ответственность между сервисами авторизации, профилей, мероприятий, геолокации, сообществ, файлового хранилища и шлюза API. Такой подход соответствует требованиям масштабируемости, независимого развития модулей и упрощения сопровождения.

\section{Цель и задачи работы}

Целью работы является проектирование и реализация серверной и инфраструктурной части платформы Underfish, обеспечивающей публикацию, поиск и сопровождение локальных мероприятий с учетом требований безопасности, масштабируемости, надежности и удобства интеграции с клиентскими приложениями.

Для достижения цели были поставлены следующие задачи:
\begin{enumerate}
    \item проанализировать предметную область и существующие решения для поиска мероприятий;
    \item сформировать функциональные и нефункциональные требования к backend-части;
    \item спроектировать доменную модель: пользователей, профили, сообщества, мероприятия, теги, локации, отзывы, посещения и файлы;
    \item выполнить переход от монолитного прототипа к микросервисной архитектуре;
    \item реализовать сервисы на Kotlin и Spring Boot с использованием PostgreSQL, JPA/Hibernate и Flyway;
    \item спроектировать механизм авторизации на основе JWT и централизованного API Gateway;
    \item подготовить Docker-инфраструктуру для локального и серверного развертывания;
    \item определить состав unit-, интеграционных и end-to-end-тестов;
    \item подготовить UML, BPMN и ER-диаграммы для описания архитектуры и бизнес-процессов.
\end{enumerate}

Объектом исследования является программная платформа поиска локальных мероприятий. Предметом исследования являются архитектурные решения backend-разработки, межсервисное взаимодействие, безопасность, DevOps-процессы и тестирование микросервисной системы.

\section{Анализ аналогов}

\subsection{Критерии сравнения}

Для анализа аналогов были выбраны критерии, значимые для гиперлокальной платформы:
\begin{itemize}
    \item простота публикации мероприятия малым организатором;
    \item наличие поиска по карте и району;
    \item поддержка сообществ и подписок;
    \item возможность фильтрации по тегам, датам и типам событий;
    \item доступность API и возможность интеграции;
    \item применимость для некоммерческих и андеграунд-событий;
    \item масштабируемость и устойчивость backend-архитектуры.
\end{itemize}

\subsection{Сравнение с существующими сервисами}

\begin{longtable}{|p{3.2cm}|p{4.2cm}|p{4.2cm}|p{3.2cm}|}
\caption{Сравнение Underfish с аналогами}\label{tab:analogs}\\
\hline
\textbf{Сервис} & \textbf{Сильные стороны} & \textbf{Ограничения} & \textbf{Вывод для Underfish} \\
\hline
\endfirsthead
\hline
\textbf{Сервис} & \textbf{Сильные стороны} & \textbf{Ограничения} & \textbf{Вывод для Underfish} \\
\hline
\endhead
Яндекс Афиша \cite{yandexafisha} & Большая аудитория, развитый поиск, узнаваемость бренда & Фокус на коммерческих событиях, высокая конкуренция за внимание пользователя, недостаточная локальность & Сделать акцент на малые события, карту района и подписки на локальные сообщества \\
\hline
Timepad \cite{timepad} & Инструменты регистрации участников, билеты, страницы мероприятий & Сервис ориентирован на формальные мероприятия и организаторов, которым нужна регистрация и продажи & Упростить публикацию для неформальных мероприятий, добавить социальный слой сообществ \\
\hline
Социальные сети & Быстрое создание постов, подписчики уже находятся в экосистеме & События теряются в ленте, слабая геофильтрация, нет единой карты мероприятий & Разделить информационный поток и каталог событий, добавить фильтры и геопоиск \\
\hline
Карты и справочники & Привязка к местам, привычный интерфейс карты & Плохо подходят для временных событий и сообществ & Использовать карту как основной способ поиска, но хранить события как отдельные сущности \\
\hline
\end{longtable}

Проведенный анализ показывает, что ниша локальных и некоммерческих событий не полностью покрыта существующими решениями. Основное отличие Underfish заключается в сочетании трех элементов: гиперлокального поиска, сообществ организаторов и простой публикации событий без необходимости создавать сложную билетную инфраструктуру.

\section{Проектирование}

\subsection{Функциональные требования}

Платформа должна поддерживать три основные группы пользователей: рядовой пользователь, организатор сообщества и администратор платформы.

Рядовой пользователь должен иметь возможность регистрироваться по email, проходить авторизацию с паролем и дополнительной проверкой, редактировать профиль, искать мероприятия на карте, фильтровать события по датам и тегам, отмечать участие, оставлять отзывы, подписываться на сообщества, получать напоминания и экспортировать событие в календарный формат ICS.

Организатор сообщества должен иметь возможность создавать и редактировать сообщество, публиковать мероприятия, прикреплять место, описание, теги и изображения, отменять или изменять мероприятие, а также инициировать уведомления подписчиков при изменении времени или места.

Администратор платформы должен иметь средства модерации: удаление пользователей, сообществ и мероприятий, а также просмотр действий, требующих аудита.

\subsection{Нефункциональные требования}

Ключевые нефункциональные требования представлены в таблице \ref{tab:nfr}.

\begin{longtable}{|p{4cm}|p{11cm}|}
\caption{Нефункциональные требования}\label{tab:nfr}\\
\hline
\textbf{Категория} & \textbf{Требование} \\
\hline
\endfirsthead
\hline
\textbf{Категория} & \textbf{Требование} \\
\hline
\endhead
Производительность & Целевое время ответа пользовательских операций не более 2 секунд при штатной нагрузке; поддержка не менее 1000 активных пользователей при горизонтальном масштабировании сервисов. \\
\hline
Безопасность & HTTPS для внешнего трафика, JWT для аутентификации, разграничение ролей, хэширование паролей, защита от SQL-инъекций за счет ORM и параметризованных запросов, фильтрация входных данных. \\
\hline
Надежность & Логирование пользовательских действий, миграции базы данных через Flyway, health-check endpoint'ы, централизованный сбор логов и метрик. \\
\hline
Масштабируемость & Разделение на сервисы, независимые базы данных, контейнеризация Docker, возможность дальнейшего переноса в Kubernetes. \\
\hline
Удобство сопровождения & Единый стиль Kotlin/Spring Boot, REST API, OpenAPI/Postman-коллекции, автоматизированные тесты и среда локального запуска. \\
\hline
\end{longtable}

\subsection{Этап 1: монолитный прототип}

Первоначально система проектировалась как монолитное приложение. На этом этапе все основные доменные модули находились в одном backend-приложении: пользователи, мероприятия, сообщества, локации, отзывы и файлы. Такой подход позволил быстро проверить требования и определить границы предметной области.

Преимущества монолитного этапа:
\begin{itemize}
    \item быстрый старт разработки и меньше инфраструктурных затрат;
    \item проще отлаживать бизнес-логику на раннем этапе;
    \item единая база данных упрощает транзакции между сущностями;
    \item удобно формировать первичную ER-модель и REST API.
\end{itemize}

Ограничения монолитного этапа:
\begin{itemize}
    \item рост связности между модулями;
    \item сложность независимого масштабирования поиска, файлов и авторизации;
    \item риск регресса при изменении одной функциональной области;
    \item затрудненное разделение ответственности в команде.
\end{itemize}

На основании этих ограничений был выбран второй этап развития -- микросервисная архитектура.

\subsection{Этап 2: микросервисная архитектура}

Микросервисный вариант разделяет систему на самостоятельные сервисы. Каждый сервис отвечает за отдельную функциональную область, имеет собственные настройки, миграции и базу данных. Взаимодействие выполняется через REST API, а для уведомлений и фоновых операций предусмотрена асинхронная модель на основе очередей сообщений \cite{kafka}.

Диаграмма компонентов подготовлена в формате Mermaid и используется в виде экспортированного PNG-изображения \texttt{uml\_component\_architecture.png}. Исходный файл диаграммы хранится в \texttt{docs/diagrams/uml\_microservices.mmd}.

\reportfigure{uml_component_architecture.png}{UML-диаграмма компонентов микросервисной архитектуры Underfish}

Основные компоненты архитектуры:
\begin{itemize}
    \item \textbf{API Gateway} -- единая точка входа клиентских приложений; маршрутизация запросов, проверка JWT, прокидывание внутренних заголовков пользователя;
    \item \textbf{Auth Service} -- регистрация, авторизация, выпуск access/refresh JWT, JWKS endpoint для проверки подписи токенов;
    \item \textbf{Profile Service} -- хранение и изменение пользовательских профилей, ролей и настроек;
    \item \textbf{Community Service} -- создание, обновление, удаление и поиск сообществ, управление участниками и организаторами;
    \item \textbf{Event Service} -- создание, редактирование, поиск и удаление мероприятий, теги, отзывы, отметки посещения;
    \item \textbf{Location Service} -- хранение координат, адресов и данных мест проведения;
    \item \textbf{File Storage Service} -- загрузка и хранение изображений мероприятий, аватаров пользователей и сообществ в S3-совместимом хранилище MinIO \cite{minio};
    \item \textbf{Notification Service} -- проектируемый сервис email-уведомлений и напоминаний о событиях.
\end{itemize}

\subsection{Модель данных}

На уровне базы данных выделяются следующие ключевые сущности:
\begin{itemize}
    \item \textbf{User} -- пользователь системы, email, хэш пароля, роль, статус учетной записи;
    \item \textbf{Profile} -- публичные данные пользователя: имя, аватар, описание, настройки уведомлений;
    \item \textbf{Community} -- локальное сообщество с названием, описанием, приватностью и статусом;
    \item \textbf{CommunityMember} -- связь пользователя с сообществом;
    \item \textbf{CommunityOrganizer} -- связь организатора с сообществом и права управления;
    \item \textbf{Event} -- мероприятие: название, описание, время начала, статус, цена, формат, организатор, сообщество и место;
    \item \textbf{Tag} -- фиксированный тег мероприятия;
    \item \textbf{EventAttendance} -- отметка пользователя «пойду»;
    \item \textbf{Review} -- отзыв пользователя после посещения;
    \item \textbf{Location} -- адрес и координаты.
\end{itemize}

\reportfigure{er_database_schema.png}{ER-диаграмма базы данных платформы Underfish}

\subsection{Бизнес-процессы}

Бизнес-процессы представлены диаграммами BPMN-style, подготовленными в Mermaid и экспортированными в PNG. Диаграммы фиксируют взаимодействие пользователя, frontend-части и backend-сервисов в ключевых сценариях платформы.

\reportfigure{bpmn_registration_login.png}{BPMN-диаграмма процесса регистрации и авторизации пользователя}

Процесс регистрации и авторизации включает следующие шаги:
\begin{enumerate}
    \item пользователь отправляет email и пароль в Auth Service;
    \item сервис проверяет уникальность email и сохраняет хэш пароля;
    \item Auth Service вызывает Profile Service для создания профиля;
    \item после успешной авторизации выпускаются access token и refresh token;
    \item клиент передает access token в API Gateway;
    \item Gateway проверяет подпись JWT через JWKS и добавляет внутренние заголовки пользователя;
    \item целевой сервис выполняет бизнес-операцию с учетом роли пользователя.
\end{enumerate}

\reportfigure{bpmn_create_event.png}{BPMN-диаграмма процесса создания мероприятия}

Процесс создания мероприятия состоит из подготовки карточки события, проверки прав организатора, сохранения основной информации, привязки геолокации, загрузки изображений и последующей публикации. При изменении времени или места события система должна формировать событие для Notification Service, чтобы подписчики и участники получили уведомления.

\reportfigure{bpmn_search_event.png}{BPMN-диаграмма процесса поиска мероприятия на карте}

\reportfigure{bpmn_join_event_review.png}{BPMN-диаграмма процесса отметки участия и публикации отзыва}

\reportfigure{uml_sequence_gateway_auth.png}{UML-диаграмма последовательности проверки авторизации через API Gateway}

\subsection{Проектирование API}

REST API проектировался вокруг ресурсов: \texttt{/api/v1/events}, \texttt{/api/v1/communities}, \texttt{/api/v1/users}, \texttt{/api/v1/files}, \texttt{/api/v1/locations}. Основные операции соответствуют HTTP-методам: \texttt{GET} для чтения, \texttt{POST} для создания, \texttt{PUT} для изменения и \texttt{DELETE} для удаления.

Для мероприятий предусмотрены фильтры по названию, статусу, онлайн-формату, диапазону дат, цене, месту, тегам, сообществу и организатору. Для сообществ предусмотрены фильтры по названию, приватности, пагинации и сортировке. Такая структура API упрощает интеграцию с frontend-приложением, Swagger/OpenAPI-документацией и Postman-коллекциями.

\section{Реализация}

\subsection{Технологический стек}

Backend реализован на языке Kotlin с использованием Spring Boot \cite{springboot}. В качестве ORM применяется Spring Data JPA/Hibernate, миграции выполняются через Flyway, основная СУБД -- PostgreSQL \cite{postgresql}. Для безопасности используется Spring Security \cite{springsecurity}. Контейнеризация выполняется средствами Docker \cite{docker}; дальнейшее масштабирование возможно через Kubernetes \cite{kubernetes}.

\begin{longtable}{|p{4cm}|p{11cm}|}
\caption{Используемые технологии}\label{tab:stack}\\
\hline
\textbf{Область} & \textbf{Технологии} \\
\hline
\endfirsthead
\hline
\textbf{Область} & \textbf{Технологии} \\
\hline
\endhead
Язык и фреймворк & Kotlin, Spring Boot, Spring Web, Spring Security, Spring Validation \\
\hline
Доступ к данным & Spring Data JPA, Hibernate, Flyway, PostgreSQL \\
\hline
Безопасность & JWT, JWKS, BCrypt/хэширование паролей, ролевой доступ, внутренний gateway-token \\
\hline
Инфраструктура & Docker, Docker Compose, VPS, отдельные контейнеры сервисов и баз данных \\
\hline
Файлы & MinIO как S3-совместимое объектное хранилище \\
\hline
Наблюдаемость & Spring Boot Actuator, Prometheus, Loki, структурированные логи \cite{prometheus,loki} \\
\hline
Тестирование & JUnit 5, Spring Boot Test, Mockito, Testcontainers/локальные профили, Postman, E2E-скрипты \\
\hline
\end{longtable}

\subsection{Реализованные сервисы}

\subsubsection{API Gateway}

Gateway является единой точкой входа для клиента. Он принимает внешние запросы, проверяет JWT, маршрутизирует запросы в нужный сервис и передает внутренние заголовки \texttt{X-User-Id}, \texttt{X-User-Roles}, \texttt{X-Internal-Token}. Такой подход уменьшает дублирование логики проверки токена в бизнес-сервисах и позволяет централизованно управлять маршрутизацией.

\subsubsection{Auth Service}

Auth Service отвечает за регистрацию, авторизацию, обновление токенов и публикацию JWKS. При регистрации сервис сохраняет учетную запись и инициирует создание профиля в Profile Service. Для авторизованных запросов используется JWT, соответствующий стандарту RFC 7519 \cite{jwt}. В перспективе сервис может быть расширен до полноценного OAuth 2.0 authorization server \cite{oauth2}.

\subsubsection{Profile Service}

Profile Service хранит пользовательские данные, роли и настройки профиля. Сервис изолирован от Auth Service, чтобы учетные данные и публичный профиль развивались независимо. Например, изменение аватара или имени не требует изменения механизма авторизации.

\subsubsection{Community Service}

Community Service реализует создание и управление сообществами. Пользователь может создать сообщество и стать организатором. В сервисе выделены операции для участников и организаторов, что позволяет отдельно контролировать подписки, права управления и модерацию.

\subsubsection{Event Service}

Event Service является центральным доменным сервисом платформы. Он обеспечивает создание и редактирование мероприятий, поиск и фильтрацию, работу с тегами, отзывами и отметками участия. Сервис взаимодействует с профилями и локациями, но хранит собственную модель события и независимую базу данных.

\subsubsection{Location Service}

Location Service отвечает за хранение адресов и координат. В дальнейшем он может быть расширен интеграцией с API Яндекс Карт или другим геокодером. Выделение геосервиса позволяет независимо оптимизировать радиусный поиск и кэширование популярных районов.

\subsubsection{File Storage Service}

File Storage Service инкапсулирует работу с объектным хранилищем. Клиент и доменные сервисы не должны напрямую зависеть от MinIO или S3 API: они получают ссылку или идентификатор файла через внутренний REST API.

\subsubsection{Notification Service}

Notification Service описан как проектируемый компонент. Он должен получать события из Kafka или другого брокера сообщений: создание мероприятия, изменение времени, отмена события, напоминание перед началом. Асинхронная отправка уведомлений снижает задержку пользовательских операций и повышает отказоустойчивость.


\subsection{Макеты и экраны реализации}

Макеты Figma и скриншоты реализованного интерфейса вынесены в отдельный иллюстративный блок. Они подтверждают, что backend API проектировался под конкретные пользовательские сценарии: регистрацию, поиск по карте, карточку мероприятия, профиль сообщества, кабинет организатора и модерацию.

Иллюстративные материалы сгруппированы в два обзорных изображения: коллаж макетов Figma и коллаж реализованных экранов. В них отражены страницы, связанные с основными сценариями backend-части: карта с маркерами мероприятий, страница поиска/каталога с фильтрами, карточка мероприятия, профиль пользователя или сообщества, экран организатора для создания мероприятия. Ключевыми экранами являются карта, поиск и карточка мероприятия, так как они напрямую связаны с Event Service и Location Service.

\begin{longtable}{|p{4.2cm}|p{5.8cm}|p{5cm}|}
\caption{Скриншоты макетов и реализации}\label{tab:mockups}\\
\hline
\textbf{Экран} & \textbf{Содержание} & \textbf{Назначение} \\
\hline
\endfirsthead
\hline
\textbf{Экран} & \textbf{Содержание} & \textbf{Назначение} \\
\hline
\endhead
Карта мероприятий & Маркеры событий, район, выбранные фильтры или радиус поиска & Подтверждает геосценарий и связь Event Service с Location Service \\
\hline
Поиск/каталог & Фильтры по дате, тегам, подпискам, цене или формату & Показывает, какие параметры должен поддерживать REST API поиска \\
\hline
Карточка мероприятия & Название, дата, место, теги, кнопка «пойду», фото & Связывает Event Service, File Storage Service и сценарий участия \\
\hline
Профиль пользователя или сообщества & Аватар, описание, подписки, опубликованные мероприятия & Показывает роль Profile Service и Community Service \\
\hline
Кабинет организатора & Форма создания/редактирования мероприятия и загрузка изображений & Обосновывает операции создания события, валидации места и загрузки файлов \\
\hline
\end{longtable}

Файл \texttt{figma\_mockups\_overview.png} представляет собой коллаж из 3--5 макетов Figma. Файл \texttt{implementation\_screens\_overview.png} оформлен аналогично, но содержит скриншоты работающего интерфейса. Макет и реализация расположены в одинаковом порядке: карта, поиск, карточка, профиль, кабинет организатора. Такое сопоставление показывает соответствие реализации первоначальному проектированию.

\reportfigure{figma_mockups_overview.png}{Обзор макетов Figma: карта, карточка мероприятия, профиль пользователя и экран организатора}

\reportfigure{implementation_screens_overview.png}{Скриншоты реализованного интерфейса и интеграции с backend API}

\subsection{Безопасность и управление входом}

Архитектура безопасности строится вокруг централизованной авторизации. Клиент проходит регистрацию или вход в Auth Service и получает пару токенов. Access token используется для короткоживущих запросов, refresh token -- для обновления сессии без повторного ввода пароля.

Основной поток авторизации:
\begin{enumerate}
    \item пользователь вводит email и пароль;
    \item Auth Service проверяет пароль по хэшу и статус учетной записи;
    \item Auth Service выпускает JWT с идентификатором пользователя и ролями;
    \item Gateway извлекает токен из заголовка \texttt{Authorization};
    \item Gateway проверяет подпись по публичному ключу из JWKS;
    \item Gateway передает во внутренний сервис только проверенную идентичность пользователя;
    \item бизнес-сервис применяет ролевые правила и проверяет владение ресурсом.
\end{enumerate}

Такой подход снижает риск подмены пользователя: внутренние сервисы доверяют не внешним заголовкам клиента, а заголовкам, сформированным Gateway после проверки токена. Дополнительно используется внутренний токен между Gateway и сервисами, чтобы прямые внешние вызовы бизнес-сервисов не считались доверенными.

\subsection{DevOps и развертывание}

Локальное окружение построено на Docker Compose. Для каждого сервиса поднимается отдельный контейнер приложения и отдельный контейнер PostgreSQL. Gateway публикуется на внешнем порту, остальные сервисы могут быть доступны через внутреннюю сеть Docker. Такой подход приближен к production-среде и удобен для интеграционного тестирования.

Основные принципы DevOps-контура:
\begin{itemize}
    \item сборка каждого сервиса из единого runtime Dockerfile с параметром имени сервиса;
    \item хранение конфигурации через переменные окружения;
    \item разделение dev, test и prod профилей Spring;
    \item независимые миграции Flyway для каждого сервиса;
    \item health-check баз данных и зависимостей;
    \item возможность развертывания на VPS и последующего переноса в Kubernetes;
    \item подготовка метрик Actuator для Prometheus и логов для Loki.
\end{itemize}

\subsection{Мониторинг и журналирование}

Для эксплуатационной надежности предусмотрены:
\begin{itemize}
    \item endpoint'ы \texttt{/actuator/health} и \texttt{/actuator/info};
    \item сбор технических метрик через Micrometer и Prometheus;
    \item централизованный сбор логов в Loki;
    \item логирование ошибок бизнес-операций и внешних интеграций;
    \item аудит действий, связанных с безопасностью: регистрация, вход, изменение ролей, удаление ресурсов.
\end{itemize}

\section{Тестирование}

\subsection{Стратегия тестирования}

Тестирование проекта разделено на несколько уровней:
\begin{enumerate}
    \item \textbf{Unit-тесты} -- проверка сервисного слоя, валидаторов, мапперов и бизнес-правил без поднятия полной инфраструктуры;
    \item \textbf{Интеграционные тесты} -- проверка Spring-контекста, репозиториев, миграций Flyway и REST-контроллеров;
    \item \textbf{Контрактное тестирование API} -- проверка соответствия Postman/OpenAPI спецификации фактическим endpoint'ам;
    \item \textbf{End-to-end-тесты} -- сценарии через Gateway: регистрация, вход, создание сообщества, создание мероприятия, поиск и отметка участия;
    \item \textbf{Нагрузочные проверки} -- оценка времени ответа поиска и создания мероприятий при росте количества пользователей.
\end{enumerate}

Целевой показатель покрытия unit-тестами составляет не менее 67\% кода сервисного слоя. Для контроллеров приоритет отдается интеграционным тестам, так как важно проверить маршрутизацию, сериализацию DTO, валидацию и безопасность.

\subsection{Тестовые сценарии}

\begin{longtable}{|p{4cm}|p{5.5cm}|p{5.5cm}|}
\caption{Основные тестовые сценарии}\label{tab:tests}\\
\hline
\textbf{Сценарий} & \textbf{Шаги} & \textbf{Ожидаемый результат} \\
\hline
\endfirsthead
\hline
\textbf{Сценарий} & \textbf{Шаги} & \textbf{Ожидаемый результат} \\
\hline
\endhead
Регистрация пользователя & Отправить email и пароль; проверить создание учетной записи; проверить создание профиля & Возвращается успешный ответ, пароль хранится в виде хэша, профиль доступен через Profile Service \\
\hline
Авторизация & Отправить корректные учетные данные; получить JWT; выполнить запрос через Gateway & Gateway пропускает запрос, целевой сервис получает идентификатор пользователя \\
\hline
Создание сообщества & Авторизованный пользователь отправляет данные сообщества & Сообщество создано, пользователь назначен организатором \\
\hline
Создание мероприятия & Организатор отправляет название, дату, место, описание, теги & Мероприятие создано, доступно в поиске, связано с сообществом и локацией \\
\hline
Поиск мероприятий & Выполнить запрос с фильтрами по датам, тегам, цене и месту & Возвращается страница мероприятий, соответствующая фильтрам \\
\hline
Отметка участия & Пользователь отмечает «пойду» & Создается запись посещения, событие появляется в списке посещаемых \\
\hline
Загрузка изображения & Отправить файл в File Storage Service & Файл сохранен в MinIO, возвращается ссылка/идентификатор \\
\hline
Удаление ресурса & Организатор или администратор удаляет событие/сообщество & Ресурс удален или помечен как удаленный, доступ запрещен неавторизованным пользователям \\
\hline
\end{longtable}

\subsection{Проверка безопасности}

Для проверки безопасности используются негативные сценарии:
\begin{itemize}
    \item запрос без JWT к защищенному endpoint'у возвращает \texttt{401 Unauthorized};
    \item запрос пользователя без роли организатора возвращает \texttt{403 Forbidden};
    \item попытка изменить чужое сообщество или мероприятие блокируется;
    \item некорректный или просроченный JWT отклоняется Gateway;
    \item прямой вызов внутреннего сервиса без \texttt{X-Internal-Token} отклоняется;
    \item SQL-инъекции в строковых фильтрах не приводят к выполнению произвольного SQL.
\end{itemize}

\subsection{Результаты тестирования}

В ходе тестирования проверены основные backend-сценарии: запуск Spring-контекста сервисов, работа контроллеров мероприятий и сообществ, миграции баз данных, маршрутизация через Gateway и базовые end-to-end-сценарии. Результаты ручной и автоматизированной проверки зафиксированы в Swagger/OpenAPI-страницах, выводе Gradle-тестов и логах Docker Compose.

\reportfigure{gradle_tests_screenshot.png}{Скриншот результатов автоматизированного тестирования}

\section{Заключение}

В ходе работы была спроектирована и реализована backend- и инфраструктурная часть платформы Underfish для поиска и публикации локальных мероприятий. Проект прошел два архитектурных этапа: монолитный прототип и микросервисную декомпозицию. Монолит позволил быстро проверить основные бизнес-сценарии, а микросервисная архитектура обеспечила разделение ответственности, независимое масштабирование и удобство сопровождения.

Были выделены сервисы авторизации, профилей, сообществ, мероприятий, геолокации, файлового хранилища, уведомлений и API Gateway. Для хранения данных применяется PostgreSQL, для файлов -- S3-совместимое хранилище MinIO, для будущих асинхронных уведомлений -- Kafka. Инфраструктура контейнеризована через Docker Compose и подготовлена к дальнейшему переносу на VPS или Kubernetes.

Особое внимание уделено безопасности: JWT, JWKS, централизованная проверка токенов на Gateway, внутренние заголовки идентичности и ролевой доступ. Также определена стратегия тестирования, включающая unit-, интеграционные, контрактные и end-to-end-проверки.

Дальнейшее развитие проекта предполагает завершение Notification Service, расширение геопоиска радиусными запросами, интеграцию с календарями, улучшение административной модерации, подключение полноценного CI/CD и накопление метрик производительности на реальной пользовательской нагрузке.

\newpage
\printbibliography[title={Список использованных источников}]

\newpage
\appendix
\section*{Приложение А. Иллюстративные материалы и диаграммы}
\addcontentsline{toc}{section}{Приложение А. Иллюстративные материалы и диаграммы}

Иллюстративные материалы проекта хранятся в каталоге \texttt{docs/diagrams}. В работе используются следующие исходники диаграмм и экспортированные изображения:

\begin{longtable}{|p{6cm}|p{9cm}|}
\caption{Перечень иллюстративных материалов}\label{tab:figures}\\
\hline
\textbf{Файл} & \textbf{Содержание} \\
\hline
\endfirsthead
\hline
\textbf{Файл} & \textbf{Содержание} \\
\hline
\endhead
\texttt{uml\_microservices.mmd} & Исходник Mermaid для UML-диаграммы микросервисов \\
\hline
\texttt{uml\_component\_architecture.png} & UML-диаграмма компонентов микросервисной архитектуры \\
\hline
\texttt{er\_database\_schema.png} & ER-диаграмма базы данных, экспортированная из инструмента проектирования БД \\
\hline
\texttt{bpmn\_registration\_login.mmd} & Исходник Mermaid для BPMN регистрации и авторизации \\
\hline
\texttt{bpmn\_registration\_login.png} & BPMN регистрации и авторизации \\
\hline
\texttt{bpmn\_create\_event.mmd} & Исходник Mermaid для BPMN создания мероприятия организатором \\
\hline
\texttt{bpmn\_create\_event.png} & BPMN создания мероприятия организатором \\
\hline
\texttt{bpmn\_search\_event.mmd} & Исходник Mermaid для BPMN поиска мероприятия пользователем \\
\hline
\texttt{bpmn\_search\_event.png} & BPMN поиска мероприятия пользователем \\
\hline
\texttt{bpmn\_join\_event\_review.mmd} & Исходник Mermaid для BPMN отметки участия и отзыва \\
\hline
\texttt{bpmn\_join\_event\_review.png} & BPMN отметки участия и публикации отзыва \\
\hline
\texttt{uml\_gateway\_auth\_sequence.mmd} & Исходник Mermaid для UML sequence diagram авторизации \\
\hline
\texttt{uml\_sequence\_gateway\_auth.png} & UML sequence diagram проверки JWT через Gateway \\
\hline
\texttt{figma\_mockups\_overview.png} & Скриншот или коллаж макетов Figma \\
\hline
\texttt{implementation\_screens\_overview.png} & Скриншоты реализованного интерфейса \\
\hline
\texttt{gradle\_tests\_screenshot.png} & Скриншот результатов тестов \\
\hline
\end{longtable}

Диаграммы в формате Mermaid экспортируются в PNG и вставляются в основной текст через команду \texttt{\textbackslash includegraphics}. Такой формат обеспечивает одинаковое отображение изображений при локальной сборке и в облачных LaTeX-редакторах.

\section*{Приложение Б. Содержание UML/BPMN-диаграмм}
\addcontentsline{toc}{section}{Приложение Б. Содержание UML/BPMN-диаграмм}

\begin{enumerate}
    \item \textbf{UML компонентов}: показать клиентское приложение, Gateway, основные сервисы, обобщенный слой хранения, Kafka, внешние интеграции и наблюдаемость без детализации каждой отдельной БД.
    \item \textbf{UML последовательности авторизации}: Client, Gateway, JWKS/Auth Service, Profile/Event/Community Service; обязательные сообщения -- login, issue token, validate token, add internal headers, execute request.
    \item \textbf{BPMN регистрации}: стартовое событие, ввод формы, проверка email, создание учетной записи, создание профиля, выдача токенов, ошибка валидации.
    \item \textbf{BPMN создания события}: выбор сообщества, проверка роли организатора, заполнение карточки, валидация места, загрузка файлов, публикация, уведомление подписчиков.
    \item \textbf{BPMN поиска}: открытие карты, ввод фильтров, запрос к Event Service, уточнение координат через Location Service, отображение маркеров, переход на карточку события.
    \item \textbf{BPMN участия и отзыва}: открытие карточки, отметка «пойду», экспорт в календарь, напоминание, проверка завершения мероприятия и сохранение отзыва.
    \item \textbf{Макеты и реализация}: общий скриншот Figma и отдельный коллаж реализованных экранов включают карту мероприятий, поиск/каталог, карточку мероприятия, профиль пользователя или сообщества, кабинет организатора.
\end{enumerate}

\end{document}
